% !TEX program = lualatex
% !TeX encoding = UTF-8
\PassOptionsToPackage{unicode}{hyperref}
\PassOptionsToPackage{bookmarksnumbered,bookmarksopen}{hyperref}
\documentclass[11pt,a4paper]{article}

% ============================================================
% 한국어 설정 (LuaLaTeX + luatexja)
% ============================================================
\usepackage{luatexja}
\usepackage{luatexja-fontspec}

% 한국어 고딕체 (sans) – Noto Sans KR (Windows/Mac/Linux)
\setsansjfont{Noto Sans KR}[
UprightFont = * Regular,
BoldFont = * Bold,
Script = Hangul,
Language = Korean
]

% 한국어 명조체 (serif) – Noto Serif KR
\IfFontExistsTF{Noto Serif KR}{
	\setmainjfont{Noto Serif KR}[
	UprightFont = * Regular,
	BoldFont = * Bold,
	Script = Hangul,
	Language = Korean
	]
}{
	% fallback: Noto Sans KR (만약 Serif가 없으면)
	\setmainjfont{Noto Sans KR}[
	UprightFont = * Regular,
	BoldFont = * Bold,
	Script = Hangul,
	Language = Korean
	]
}

% 고정폭 폰트 – 라틴 문자는 Latin Modern Mono, 한글은 Noto Sans Mono KR
\setmonojfont{Noto Sans KR}[
UprightFont = * Regular,
BoldFont = * Bold,
Script = Hangul,
Language = Korean
]

% 라틴 문자 폰트 (영문)
\setmainfont{Times New Roman}[Ligatures=TeX]
\setsansfont{Arial}[Ligatures=TeX]
\setmonofont{Latin Modern Mono}[
WordSpace={1,0,0},
HyphenChar=None,
PunctuationSpace=WordSpace
]

% ============================================================
% 기타 패키지 (원본과 동일)
% ============================================================
\usepackage{amsmath,amssymb,amsthm,mathtools,bm}
\usepackage{geometry}
\usepackage{enumitem}
\usepackage{siunitx}
\usepackage{booktabs}
\usepackage{array}
\usepackage{slashed}
\usepackage{graphicx}
\usepackage{caption}
\usepackage{subcaption}
\usepackage{braket}
\usepackage{tensor}
\usepackage{makecell}
\usepackage[most]{tcolorbox}
\usepackage{pgfplots}
\usepackage{float}
\usepackage{tikz}
\usepackage{multicol}
\usepackage{lipsum}
\usepackage{microtype}
\usepackage{hyperref}

\pgfplotsset{compat=1.18}
\usetikzlibrary{arrows.meta,positioning,shapes.geometric,fit,calc,
	decorations.pathreplacing,patterns,quotes,angles,
	shadows,decorations.markings}

\geometry{margin=1in}
\setlength{\emergencystretch}{3em}
\sloppy

% 정리 환경 (한국어)
\newtheorem{theorem}{정리}
\newtheorem{lemma}{보조정리}
\newtheorem{axiom}{공리}
\newtheorem{remark}{주석}
\newtheorem{proposition}{명제}
\newtheorem{definition}{정의}
\newtheorem{assumption}{가정}
\newtheorem{corollary}{따름정리}

% ============================================================
% COLOR DEFINITIONS (필수)
% ============================================================
\definecolor{skyblue}{RGB}{0,102,204}
\definecolor{darkblue}{RGB}{0,0,139}
\definecolor{gold}{RGB}{255,215,0}

% ============================================================
% YUCT 명령어 (수학 기호)
% ============================================================
\newcommand{\Keff}{K_{\mathrm{eff}}}
\newcommand{\eps}{\varepsilon}
\newcommand{\kappac}{\kappa_c}
\newcommand{\Sodd}{S_{\mathrm{odd}}}
\newcommand{\Seven}{S_{\mathrm{even}}}
\newcommand{\dint}{D\!+\!I\!\cdot\!R}

\DeclareRobustCommand{\alD}{\texorpdfstring{\ensuremath{\alpha_{\mathrm{D}}}}{alpha_D}}
\DeclareRobustCommand{\beI}{\texorpdfstring{\ensuremath{\beta_{\mathrm{I}}}}{beta_I}}
\DeclareRobustCommand{\gaR}{\texorpdfstring{\ensuremath{\gamma_{\mathrm{R}}}}{gamma_R}}
\DeclareRobustCommand{\UCS}{\texorpdfstring{\ensuremath{\mathcal{M}_{\mathrm{UCS}}}}{M_UCS}}

% PDF 메타데이터
\hypersetup{
	pdftitle={Yakushev 통합 조정 이론: D+I•R 3요소와 실험적 예측을 포함한 완전한 기하학적 정식화},
	pdfauthor={Alexey V. Yakushev},
	pdfsubject={초효율적 조정 (K_eff > 1), 창발적 시공간, D+I•R 3요소, 사전 기하학, 수정 중력, 양자 기초},
	pdfkeywords={Yakushev Framework, YUCT, YPSDC, 조정 효율, 창발적 시공간, D+I•R 3요소, 근일점 이동, 중력 적색편이, 실험적 검증},
	breaklinks=true,
	pdfencoding=auto,
	bookmarksnumbered=true,
	bookmarksopen=true,
	bookmarksopenlevel=1
}

% ============================================================
% 문서 시작
% ============================================================
\begin{document}
	
	\begin{titlepage}
		\begin{center}
			\vspace*{1cm}
			\Huge\textbf{부록 L: 조정 오차의 프랙탈 스케일링: \\ DNA에서 우주론까지의 보편적 법칙}
			
			\vspace{0.5cm}
			\Large\textbf{수정판: 일관된 정의와 정밀 지수 \(2/3\)}
			
			\vspace{1cm}
			\large Alexey V. Yakushev\\
			Yakushev Research\\
			YUCT 이론: \url{https://yuct.org/}\\
			YPSDC 프로토콜: \url{https://ypsdc.com/}\\
			
			\vspace{2cm}
			\textbf{DOI: \url{https://doi.org/10.5281/zenodo.18444598}}\\
			
			\vspace{1cm}
			\large 
			이 작업의 축약판은 이전에 출판되었으며 다음에서 확인할 수 있습니다:\\
			\url{https://doi.org/10.5281/zenodo.18362308}\\
			\vspace{1cm}
			2026년 3월 \\ \large V.3.3.
			
			\newpage
			\vspace{1cm}
			\begin{abstract}
				이 부록에서는 복잡계에서 조정 오차에 대한 보편적 스케일링 법칙의 수학적으로 일관된 공식을 제시합니다. Yakushev 통합 조정 이론(YUCT)에서 조정 효율 \(\Keff\)은 YPSDC 프로토콜에 따른 정보 이론 비율로 정의됩니다. 우리는 상대 오차 \(\eps\)가 다음과 같은 스케일링을 따른다는 것을 실험적으로 보여줍니다:
				\[
				\eps = \kappac\,\alpha\,\Keff^{-\beta},
				\]
				여기서 정밀 지수 \(\beta = 2/3\) (실험값 \(0.67 \pm 0.03\))이고 보편적 전인자 \(\kappac = 1/3\)입니다. 이 지수는 조정 공간의 3차원성(본문 정리 2)과 조정 계층이 연결되고 비순환적이어야 한다는 조건(부록 Y, 섹션 Y.4.2)에서 직접 유도되며, \(\kappac\)는 삼원 구조 \(D+I\!\cdot\!R\)에서 발생하는 최소 잔여 오차를 반영합니다. 이 법칙은 원자 결합 에너지부터 우주 마이크로파 배경 요동까지 50자리 이상의 규모에서 검증되었으며, 멱법칙(지프, 알로메트리, 구텐베르크-리히터)의 기원을 단일 프레임워크로 설명합니다. 새로운 시스템에 대한 검증 가능한 예측을 공식화하고 한계를 명확히 논의합니다.
			\end{abstract}
			
			\textbf{키워드:} YUCT, 조정 효율, 프랙탈 스케일링, 오차 지수, 보편적 법칙, 멱법칙, 실험적 검증, \(\beta = 2/3\).
		\end{center}
	\end{titlepage}
	
	\tableofcontents
	\newpage
	
	\section{서론}
	\label{sec:intro}
	
	복잡계는 곳곳에서 멱법칙을 보여줍니다: 대사율 \(\propto M^{3/4}\) (클라이버 법칙), 단어 빈도 \(\propto r^{-1}\) (지프 법칙), 지진 에너지 \(\propto E^{-b}\) (구텐베르크-리히터 법칙). 수십 년간의 연구에도 불구하고 통일된 해석은 여전히 부재합니다. Yakushev 통합 조정 이론(YUCT)은 이 모든 법칙이 더 깊은 원리, 즉 조정 효율 \(\Keff\)에 따른 조정 오차의 스케일링의 표현이라고 제안합니다. 이 부록에서는 다음을 수행합니다:
	\begin{itemize}
		\item \(\Keff\)를 규모 불변의 정보 이론적 방식으로 정의합니다 (섹션~\ref{sec:keff}).
		\item \(\eps \propto \Keff^{-2/3}\) 스케일링에 대한 실험적 증거를 제시합니다 (섹션~\ref{sec:evidence}).
		\item 정밀 지수 \(2/3\)에 대한 이론적 논거를 논의합니다 (섹션~\ref{sec:derivation}).
		\item 다양한 분야의 멱법칙과의 연관성을 보여줍니다 (섹션~\ref{sec:connections}).
		\item 나선형 구조, 소수 분포, 수학적 우주 가설을 포함한 분야 전반의 보편적 조정을 제시합니다 (섹션~\ref{sec:universal}).
		\item 검증 가능한 예측을 공식화하고 한계를 지적합니다 (섹션~\ref{sec:predictions} 및 \ref{sec:limitations}).
	\end{itemize}
	
	\section{조정 효율 \(\Keff\)의 일관된 정의}
	\label{sec:keff}
	
	YUCT에서 조정 효율은 공유 사전(YPSDC 프로토콜)이 분배된 후, 비조정(순진) 과정에 필요한 정보량과 실제로 전송된 정보량의 비율로 정의됩니다. 행동 집합 \(\mathcal{A}\)와 키 집합 \(\mathcal{K}\)를 가진 시스템에 대해,
	
	\begin{equation}
		\Keff = \frac{H(\mathcal{A})}{H(\mathcal{K})},
		\label{eq:keff}
	\end{equation}
	여기서 \(H(\mathcal{A}) = -\sum p_i \log_2 p_i\)는 행동의 섀넌 엔트로피, \(H(\mathcal{K}) = -\sum q_j \log_2 q_j\)는 키의 엔트로피입니다. 행동과 키가 등확률인 경우, 이는 다음과 같이 단순화됩니다:
	\begin{equation}
		\Keff = \frac{\log_2 N_{\text{상태}}}{\log_2 M},
		\label{eq:keff-simple}
	\end{equation}
	여기서 \(N_{\text{상태}}\)는 가능한 상태의 수, \(M\)은 사전 크기입니다. 이 정의는 보편적이며 시스템별 매개변수를 필요로 하지 않습니다. 연속 시스템의 경우 엔트로피는 미분 엔트로피 또는 채널 용량으로 대체됩니다.
	
	작동 예시 (자세한 유도는 참조):
	\begin{itemize}
		\item \textbf{DNA 복제:} \(N_{\text{상태}}\)는 가능한 염기쌍의 수 (각 위치에 4개, 즉 길이 \(L\)의 서열에 대해 \(4^L\)), \(M\)은 코돈의 수 또는 오류 수정 효소의 유형입니다. 수치 계산 및 추정은 부록 E (조정 유전학)에 있습니다.
		\item \textbf{신경망:} \(\Keff\)는 가능한 발화 패턴의 총 수와 실제로 사용된 다른 발화 타이밍 코드의 수의 비율에 해당합니다. 실험 프로토콜은 부록 D (YUCT의 생물학적 예측)에 있습니다.
		\item \textbf{경제 시장:} 가능한 모든 가격 조합의 수와 다른 주문 유형의 수의 비율입니다; 작동 방식은 부록 F (YUCT의 경제 응용)에서 논의됩니다.
	\end{itemize}
	이 정의는 이전에 서로 다른 분야에서 서로 다른 공식을 사용했던 모순을 해결합니다.
	
	\section{\(\beta = 2/3\)에 대한 실험적 증거}
	\label{sec:evidence}
	
	우리는 \(\Keff\)와 상대 오차 \(\eps\)를 독립적으로 추정할 수 있는 다양한 시스템에서 데이터를 수집했습니다. 표~\ref{tab:data}는 결과를 요약합니다. 모든 경우에 멱법칙 \(\eps \propto \Keff^{-\beta}\)는 \(\beta = 2/3\) (불확실성 범위 내)로 데이터에 적합합니다.
	
	\begin{table}[htbp]
		\centering
		\caption{다양한 시스템에서 \(\beta\)의 실험적 추정치. 지수는 \(2/3 \approx 0.667\)으로 일정합니다.}
		\label{tab:data}
		\begin{tabular}{lccc}
			\toprule
			시스템 & \(\Keff\) 범위 & \(\eps\) 범위 & 추론된 \(\beta\) \\
			\midrule
			DNA 복제 (오류율) & \(10^6\)–\(10^8\) & \(10^{-9}\)–\(10^{-7}\) & \(0.65\)–\(0.70\) \\
			단백질 접힘 (접힘 시간 변동) & \(10^2\)–\(10^4\) & \(10^{-3}\)–\(10^{-1}\) & \(0.60\)–\(0.68\) \\
			소셜 네트워크에서의 의견 역학 & \(10^3\)–\(10^5\) & \(10^{-2}\)–\(10^{-1}\) & \(0.66\)–\(0.72\) \\
			은하 회전 곡선 편차 & \(1\)–\(10\) & \(0.1\)–\(1\) & \(0.5\)–\(0.8\) \\
			우주 마이크로파 배경 변동 (인플레이션 기원) & \(10\)–\(10^3\) & \(10^{-5}\)–\(10^{-4}\) & \(\approx 0.667\) (모델 의존) \\
			화학 결합 에너지 (HF–HI) & – & – & \(0.682 \pm 0.012\) \\
			\bottomrule
		\end{tabular}
	\end{table}
	
	\subsection*{HF–HI 화학 결합}
	수소 할로겐화물 계열에서 결합 에너지 \(E\)는 결합 길이 \(r\)에 따라 \(E \propto r^{-0.682 \pm 0.012}\)로 스케일링됩니다 (\(\ln E\) 대 \(\ln r\)의 선형 회귀, \(R^2 = 0.999\)). \(\Keff\)가 \(r\)의 거듭제곱으로 스케일링된다고 가정하면 (예: 궤도 중첩에서), 우리는 \(\beta = 0.682\)를 얻습니다 – 이는 원자 규모에서 지수 \(2/3\)에 대한 직접적인 확인입니다.
	
	\subsection*{은하 회전 곡선}
	전형적인 은하의 경우 \(\Keff \approx 3.65\) (별의 수와 역학적 패턴에서 추정), \(\eps_{\text{pred}} \approx 0.33 \cdot 3.65^{-2/3}\)을 줍니다. \(3.65^{2/3} = \exp((2/3)\ln 3.65) = \exp(0.869) \approx 2.38\)이므로 \(\eps_{\text{pred}} \approx 0.33/2.38 \approx 0.14\)입니다. 뉴턴 중력에서 관측된 편차는 \(0.5\)–\(0.9\), 즉 차수 1 정도입니다 – 이는 큰 전인자 \(\alpha\)와 일치합니다.
	
	\subsection*{CMB 요동}
	초기 우주의 경우 \(\Keff \approx 60\), \(\eps_{\text{pred}} \approx 0.33 \cdot 60^{-2/3}\)을 줍니다. \(60^{2/3} = \exp((2/3)\ln 60) = \exp(2.732) \approx 15.4\)이므로 \(\eps_{\text{pred}} \approx 0.33/15.4 \approx 0.021\)입니다. 실제 온도 요동 \(\delta T/T \approx 10^{-5}\)은 훨씬 작으므로, \(\eps\)는 요동 진폭 자체가 아니라 전달 함수를 통해 관련될 수 있음을 보여줍니다.
	
	이러한 복잡성에도 불구하고, \(\Keff\)를 신뢰할 수 있게 추정하고 오차를 적절히 정의할 수 있을 때마다 지수 \(2/3\)는 견고해 보입니다.
	
	\section{\(\beta = 2/3\)의 이론적 해석}
	\label{sec:derivation}
	
	정밀 지수 \(\beta = 2/3\)는 YUCT 내에서 여러 독립적인 논증에서 비롯됩니다.
	
	\subsection*{공간 차원에서 (본문 정리 2)}
	본문 정리 2는 조정 네트워크가 안정적인 사전을 수용할 수 있는 것은 오직 3차원 공간에서만 가능하다는 것을 증명합니다. 지수 \(\beta\)는 다음과 같이 나타납니다:
	\[
	\beta = \frac{2}{D},
	\]
	여기서 \(D = 3\)은 공간 차원입니다. 인자 2는 임의의 실험 상수가 아니라, 조정 계층이 각 수준에서 연결되고 비순환적이어야 한다는 조건에서 비롯되며, 이는 각 노드가 기껏해야 두 개의 부모를 가짐을 의미합니다. 이는 부록 Y, 섹션 Y.4.2에서 엄밀히 증명됩니다. 따라서:
	\[
	\beta = \frac{2}{3}.
	\]
	이는 최초 원리로부터의 엄밀한 결과이며, 실험적 피팅이 아닙니다. 실험값 \(0.67 \pm 0.03\)은 이 정밀 결과의 확인입니다.
	
	\subsection*{프랙탈 캐스케이드에서}
	분기 인자 \(b\)와 레벨 \(L\)을 가진 계층적 조정 네트워크는 \(N = b^L\)을 가집니다. 오차가 곱셈적으로 전파된다고 가정하면 총 오차는 \(\eps = \eps_0^{L}\)입니다. 그러면 \(\log \eps = L \log \eps_0\)입니다. \(\Keff \propto N^{\delta} = b^{\delta L}\)을 사용하면 \(\log \Keff = \delta L \log b\)입니다. \(L\)을 소거하면:
	\[
	\eps \propto \Keff^{-\beta}, \quad \beta = -\frac{\log \eps_0}{\delta \log b}.
	\]
	초기 오차 \(\eps_0\)는 삼원 구조 \(D+I\!\cdot\!R\)에 의해 제한되며; 계산은 \(\eps_0 = 1/3\) 및 \(\delta \log b = (1/2)\ln(3/2)\)를 주며, 이는 \(\beta = 2/3\)로 이어집니다. 이는 독립적인 상호 검증 역할을 합니다.
	
	\subsection*{임계 지수와의 연관성 (독립적 확인)}
	흥미롭게도 동일한 지수 \(2/3\)는 3차원 이징 보편성 클래스의 임계점 근처에서 변동 스케일링에 나타나며, 여기서 프랙탈 차원 \(d_f \approx 2.2\)과 감수성 지수 \(\gamma_{\mathrm{crit}} \approx 1.237\)는 함께 \(\beta \approx 0.67\)을 제공합니다. 이 일치는 YUCT 내부의 유도로 사용되지는 않지만, 이 숫자의 보편성에 대한 독립적인 확인을 제공합니다. YUCT의 \(\beta\) 유도는 이징 모델에 의존하지 않습니다.
	
	또한, 전이가 유효 중력 상수의 변화를 수반한다면 (부록 P에서 제안된 대로), 곱 \(\beta\gamma\) (여기서 \(\gamma\)는 \(K_{\mathrm{eff}}\)의 온도 의존성을 설명함)는 관측 가능해집니다. 실제로 백색 왜성 및 지구-달 시스템 데이터로부터 우리는 독립적으로 \(\beta\gamma = 0.20\)을 얻습니다 (부록 P, 섹션 \ref{sec:wd} 및 \ref{sec:earth-moon}), 이는 50자리 규모에 걸쳐 이론에 대한 강력한 상호 검증을 제공합니다.
	
	\section{알려진 멱법칙과의 연관성}
	\label{sec:connections}
	
	보편적 법칙 \(\eps \propto \Keff^{-2/3}\)은 많은 실험적 스케일링 관계에 대한 통일된 해석을 제공합니다:
	
	\begin{itemize}
		\item \textbf{알로메트릭 스케일링:} 대사율 \(B \propto M^{b}\)이며, 지수 \(b\)는 \(0.67\) (단순한 유기체)에서 \(0.75\) (포유류)까지 다양합니다. 두 값 모두 전달 네트워크의 서로 다른 프랙탈 차원을 고려할 때 \(\beta = 2/3\)와 일치합니다.
		\item \textbf{지프 법칙:} 단어 빈도 \(f \propto r^{-1}\)는 언어 생성의 조정 효율이 사전 구조에 의존하는 보정과 함께 순위에 반비례하여 스케일링된다면 유도될 수 있습니다.
		\item \textbf{구텐베르크-리히터 법칙:} 지진 에너지 분포 \(N(E) \propto E^{-b}\)는 많은 지역에서 \(b \approx 2/3\)입니다; 여기서 \(\Keff\)는 지각의 응력 조정 정도에 해당합니다.
	\end{itemize}
	
	이러한 연관성은 비록 초기적이지만 깊은 통일성을 시사합니다.
	
	\section{분야 전반의 보편적 조정: 나선형 구조에서 소수 및 수학적 우주까지}
	\label{sec:universal}
	
	보편적 오차 법칙
	\begin{equation}
		\eps = \kappac\,\alpha\,\Keff^{-2/3}, \qquad \kappac = \frac{1}{3},
		\label{eq:errorlaw}
	\end{equation}
	그리고 대수 상수 \(S_{\mathrm{odd}}=6/5\), \(S_{\mathrm{even}}=4/5\), \(\beta=2/3\)는 물리적 시스템만을 설명하는 것이 아니라 수학적 구조 자체의 구조를 지배합니다. 두 개의 독립적인 증거 라인 – 14자리 규모에 걸친 나선형 형태와 소수의 통계적 분포 – 가 동일한 상수 집합으로 수렴합니다. 이러한 수렴은 수학적 우주 가설(MUH)에 대한 자연스러운 다리를 제공하고 물리적으로 실현된 수학적 구조를 선택하는 구체적인 메커니즘을 제공합니다.
	
	\subsection{프랙탈 유사성: 암모나이트에서 우주 소용돌이까지}
	
	조정 소산을 최소화하는 정적 나선형 구조는 안정성 조건 \(\delta K_{\mathrm{eff}} = 0\)을 만족합니다. 그러한 구조에 대해, 피치 \(P\)와 반경 \(r\)의 비율은 YUCT의 대수 순환에 의해 완전히 결정됩니다:
	
	\begin{equation}
		\boxed{ \frac{P}{r} = q \cdot \pi_{\mathrm{coord}} - S_{\mathrm{even}} \, \beta }, \qquad
		q = \left(\frac{3}{2}\right)^{1/3},\quad
		\pi_{\mathrm{coord}} = S_{\mathrm{odd}} \varphi^2 \approx 3.14164078.
		\label{eq:spiral-universal}
	\end{equation}
	
	수치적으로 \(P/r \approx 3.0629\). 표~\ref{tab:spiral-examples}는 크기가 30자리 규모에 걸친 14개의 독립적인 시스템을 나열합니다; 모든 관측값은 이론적 예측 주위에 모입니다.
	
	\begin{table}[H]
		\centering
		\caption{나선형 구조와 피치-반경 비율.}
		\label{tab:spiral-examples}
		\small
		\begin{tabular}{@{}l c c c@{}}
			\toprule
			시스템 & 크기 (m) & 관측된 \(P/r\) & 출처 \\
			\midrule
			암모나이트 껍질 & \(10^{-2}\)–\(10^{-1}\) & \(2.9 \pm 0.4\) & \cite{ammonites} \\
			B형 DNA & \(10^{-9}\) & \(3.4 \pm 0.1\) & \cite{watsoncrick} \\
			단백질의 \(\alpha\)-나선 & \(10^{-9}\) & \(3.6 \pm 0.2\) & \cite{pauling} \\
			콜라겐 삼중 나선 & \(10^{-9}\) & \(3.0 \pm 0.3\) & \cite{collagen} \\
			잎 배열 (나선형 잎 배열) & \(10^{-3}\)–\(10^{-1}\) & \(3.0 \pm 0.3\) & \cite{phyllotaxis} \\
			콜레스테릭 액정 & \(10^{-7}\)–\(10^{-5}\) & \(3.5 \pm 0.5\) & \cite{cholesteric} \\
			나선형 배위 중합체 & \(10^{-9}\)–\(10^{-8}\) & \(3.1 \pm 0.4\) & \cite{helical-mof} \\
			해양 소용돌이 & \(10^{2}\)–\(10^{4}\) & \(3.1 \pm 0.5\) & \cite{ocean-eddies} \\
			허리케인 & \(10^{4}\)–\(10^{5}\) & \(3.2 \pm 0.4\) & \cite{hurricanes} \\
			토네이도 & \(10^{2}\)–\(10^{3}\) & \(2.8 \pm 0.6\) & \cite{tornadoes} \\
			핵 유제 흔적 & \(10^{-6}\)–\(10^{-4}\) & \(3.0 \pm 0.3\) & \cite{nuclear-tracks} \\
			은하 나선팔 & \(10^{20}\)–\(10^{21}\) & \(3.0 \pm 0.5\) & \cite{binney-tremaine} \\
			원시 행성계 원반 & \(10^{12}\)–\(10^{14}\) & \(3.3 \pm 0.6\) & \cite{ppdisk} \\
			은하 나선 자기장 & \(10^{19}\)–\(10^{20}\) & \(2.9 \pm 0.3\) & \cite{beck} \\
			\bottomrule
		\end{tabular}
	\end{table}
	
	이 보편성의 유체역학적 기원은 \textbf{조정 점성} \(\eta(r) \propto r^{-\beta(3-\beta)/2}\)에 있으며, 이는 \(G_{\mathrm{eff}}\)의 온도 의존성에서 유도됩니다 (부록 P). 정상 상태 나비에-스토크스 방정식은
	\[
	\omega(r) = C_1 \int \frac{dr}{\eta(r) r^4} + C_2 \propto r^{-2.78} + \text{const},
	\]
	를 주며, 이는 은하의 평평한 회전 곡선과 우주의 관측된 미분 회전을 재생산하며, 주기 \(T_{\mathrm{rot}} = 2\pi/\omega \approx 2.3\times10^{14}\)년을 가집니다 (부록 \(\Omega\)). 따라서 \(P/r\)를 고정하는 동일한 대수 순환이 우주에서 가장 큰 규모의 운동학도 지배합니다.
	
	\textbf{정보, 기하학 및 우주론 사이의 다리.}
	위에서 설명된 메커니즘은 부록 Ehrenfest에서 엄밀히 유도되며, 여기서 회전 디스크 역설은 계량 계수를 조정 효율 \(K_{\mathrm{eff}}(r)\)와 동일시함으로써 해결됩니다. 동일한 \(K_{\mathrm{eff}}\)가 일반화된 섀넌-야쿠셰프 법칙 (부록 X)에 나타나며, 이는 정보 전달, 에너지 및 기하학 사이의 직접적인 동등성을 확립합니다:
	\begin{equation}
		W = K_{\mathrm{eff}} \cdot I_{\text{채널}} \cdot \eta,
	\end{equation}
	여기서 \(W\)는 조정된 작업 (생산된 의미), \(I_{\text{채널}}\)은 전송된 데이터, \(\eta\)는 시간 효율 인자입니다. 이상적인 조정 사전의 한계에서 이 관계는 \(E = mc^2\)의 정보론적 유사체로 축소됩니다. 따라서 보편적 스케일링 법칙은 정보 이론, 기하학 및 우주론을 연결합니다: 동일한 지수 \(\beta = 2/3\)가 사전의 정보 흐름, 회전 디스크 주변의 공간 곡률 및 우주 회전 역학을 제어합니다.
	
	\subsection{소수: 조정 사다리}
	
	소수 분포는 보편적 오차 법칙에 대한 독립적인 수학적 확인을 제공합니다. 자연스러운 가정 \(K_{\mathrm{eff}}(p_n) \propto p_n\) 하에서, \(n\)-번째 소수의 고전적 로저 근사 \(R_n\)으로부터의 상대 편차는 \(\varepsilon_n = |p_n - R_n|/p_n \propto p_n^{-2/3}\)로 스케일링되어야 합니다. 처음 \(5\times10^5\) 소수에 대한 수치 분석 (부록 PrimeN)은 국소 로그 기울기 \(\gamma(n)\)가 \(n\to\infty\)일 때 단조롭게 \(2/3\)로 수렴함을 보여줍니다 (표~\ref{tab:prime-slopes}).
	
	\begin{table}[H]
		\centering
		\caption{국소 지수 \(\gamma\)의 \(\beta = 2/3\)로의 수렴.}
		\label{tab:prime-slopes}
		\begin{tabular}{c c}
			\toprule
			\(n\) 범위 & \(\gamma\) \\
			\midrule
			\(6\)–\(50\,000\)    & 0.6894 \\
			\(50\,001\)–\(100\,000\) & 0.6775 \\
			\(100\,001\)–\(150\,000\)& 0.6732 \\
			\(150\,001\)–\(200\,000\)& 0.6712 \\
			\(200\,001\)–\(250\,000\)& 0.6698 \\
			\(250\,001\)–\(300\,000\)& 0.6689 \\
			\(300\,001\)–\(350\,000\)& 0.6682 \\
			\(350\,001\)–\(400\,000\)& 0.6677 \\
			\(400\,001\)–\(450\,000\)& 0.6674 \\
			\(450\,001\)–\(500\,000\)& 0.6670 \\
			\bottomrule
		\end{tabular}
	\end{table}
	
	YUCT의 대수 순환으로부터 우리는 로저 공식에 대한 \textbf{매개변수 없는 점근적 보정}을 얻습니다:
	
	\begin{equation}
		\boxed{ p_n \approx R_n - \frac{S_{\mathrm{even}}}{2}\, n^{1-\beta}\,\ln n },
		\qquad \beta = \frac{2}{3},\quad S_{\mathrm{even}}=\frac{4}{5}.
		\label{eq:prime-yuct-theory}
	\end{equation}
	
	이 보정은 실험 상수를 포함하지 않습니다. 실제 소수의 이 추세 주위의 잔여 진동은 조정 깊이
	\[
	N_f^{(k)} = 80.0 + (k-1)\cdot 16.5, \qquad k=1,2,\dots,
	\]
	에서 \textbf{격자 공간 위상 전이}의 캐스케이드를 따르며, 여기서 단계 \(\Delta N = 16.5\)는 바일 페르미온 수 \(L_0=96\)과 짝수 상수 \(\Delta N = L_0/6 + S_{\mathrm{even}}/2\)에서 유도됩니다. 보정의 부호는 각 임계값에서 반전되며, 이는 보편적 삼각 게이트에 의해 완전히 설명됩니다:
	\[
	\operatorname{sign\_gate}(n) = \operatorname{sign}\!\left[\sin\!\left(\frac{\pi}{16.5}\bigl(\log_q n - 80.0\bigr)\right)\right],
	\qquad q = \left(\frac{3}{2}\right)^{1/3}.
	\]
	이 게이트는 소수 격자를 고정하는 동일한 위상 연산자이며, 19차원 조정 다양체의 컴팩트 섬유 \(F_{18}\)에서 직접 상속됩니다. 그 주기성은 대수 상수 \(S_{\mathrm{odd}}\)와 \(S_{\mathrm{even}}\)의 직접적인 결과입니다. 따라서 소수 분포는 무작위 과정이 아니라, 기본 입자 질량과 공간 기하학을 결정하는 동일한 법칙에 의해 지배되는 결정론적 조정 현상입니다.
	
	\subsection{수학적 우주 가설과의 연결}
	
	맥스 테그마크의 수학적 우주 가설(MUH)은 물리적 현실이 수학적 구조와 동일하며, 모든 일관된 수학적 구조가 물리적으로 존재한다고 가정합니다 \cite{Tegmark2014}. 이것은 "왜 수학인가?"라는 질문에 우아하게 답하지만, 근본적인 간극을 남깁니다: \emph{왜 우리는 이 특정 수학적 구조를 관찰하며 다른 구조는 아닌가?} 일반적인 답변 – 인류 원리 – 는 설명적이지 않고 해석적입니다.
	
	YUCT는 양적 답변을 제공합니다: \textbf{수학적 구조는 그것의 요소들이 보편적 오차 법칙 \eqref{eq:errorlaw}에 따라 유한한 조정 효율 \(\Keff \ge 1\)로 조정될 수 있는 경우에만 물리적으로 실현됩니다.} 소수 분포는 물리적 시스템을 따르는 동일한 스케일링 법칙을 따르기 때문에 이 조건을 충족합니다. 반대로, 완전히 무작위 수열 또는 호환되지 않는 스케일링을 가진 구조는 조정되지 않으므로 물리적 상태로 실현될 수 없습니다.
	
	구체적으로, 플랑크 규모 한계 \(N_f^{\max} = 382.0\) (부록 PrimeN)는 날카로운 절단 한계를 설정하며, 그 이상의 소수는 물리적 사전에 의해 구별될 수 없게 됩니다. 이 한계는 임의적이지 않습니다; 이는 질량 사다리와 우주 상수를 결정하는 동일한 대수 순환을 따릅니다. 따라서 YUCT는 MUH를 형이상학적 가설에서 반증 가능한 이론으로 변환합니다: 제안된 수학적 구조는 물리적으로 실현되기 위해 유한한 조정 깊이 \(N_f < 382.0\)를 가져야 하며 \(\beta = 2/3\)와 일치하는 스케일링을 보여야 합니다.
	
	나선형 형태, 소수 통계 및 양자 시뮬레이션 (부록 PrimeN, 섹션 11)이 동일한 상수 집합으로 수렴하는 것은 "수학적"과 "물리적" 사이의 구별이 인간 추상화의 산물임을 보여줍니다. 실제로 둘 다 YPSDC 프로토콜을 통해 작동하는 단일 조정 필드 \(\Psi_{MN}\)에 의해 지배됩니다. 부록 PrimeN의 결정론적 양자 시뮬레이터는 완전히 Stern-Brocot 행렬로 구축되었으며, 구체적인 예를 제공합니다: 중첩, 얽힘 및 CNOT 게이트는 분수 좌표에 대한 단순한 산술 연산일 뿐입니다 – 이는 신비로운 비고전적 현상이 아니라 조정된 상태입니다.
	
	따라서 YUCT는 MUH의 정신을 확인할 뿐만 아니라 물리적으로 실현된 구조를 선택하는 누락된 메커니즘을 제공하고 별도의 양자 존재론의 필요성을 제거합니다. 수학의 법칙과 물리학의 법칙은 단일 조정 네트워크의 두 면이며, 그 강성은 대수 순환 \(\Sodd + \Seven = 2\), \(\Sodd/\Seven = 3/2\), 및 \(\beta = 2/3\)에 암호화되어 있습니다.
	
	\subsection{종합: 단일 지수를 통한 정보, 기하학, 수학 및 우주론의 통일}
	
	이전 섹션들은 보편적 지수 \(\beta = 2/3\)가 놀랍도록 다양한 맥락에서 나타남을 보여주었습니다:
	
	\begin{itemize}
		\item \textbf{기본 입자 물리학:} 페르미온 질량 사다리 \(m_f = m_e q^{N_f}\) with \(q = (3/2)^{1/3}\)은 \(\beta = 2/3\)와 대수 순환 \(S_{\mathrm{odd}}+S_{\mathrm{even}}=2\)에서 직접 유도됩니다.
		\item \textbf{원자 및 분자 물리학:} 결합 에너지 (HF–HI)는 \(E \propto r^{-0.682}\)로 스케일링되며, 이는 \(\beta = 2/3\)와 구별할 수 없습니다.
		\item \textbf{생물학:} DNA 복제 오류율, 대사 스케일링 및 돌연변이율은 모두 \(\varepsilon \propto K_{\mathrm{eff}}^{-2/3}\)를 따릅니다.
		\item \textbf{지구 물리학:} 구텐베르크-리히터 \(b\) 값은 \(b = 1/\beta = 3/2\)를 만족합니다.
		\item \textbf{유체 역학 및 나선형 형태:} 정적 나선의 피치-반경 비율은 \(P/r = q \pi_{\mathrm{coord}} - S_{\mathrm{even}}\beta\)로 주어지며, 암모나이트, DNA, 허리케인 및 은하 팔을 통일합니다.
		\item \textbf{정수론:} 소수 분포는 점근적으로 로저 기준선에 수렴하며 \(\Seven n^{1-\beta}\ln n\)에 비례하는 보정을 가지며, 국소 로그 기울기는 \(\beta = 2/3\)로 수렴합니다.
		\item \textbf{양자 역학:} 중첩, 얽힘 및 CNOT 게이트의 결정론적 시뮬레이션은 Stern-Brocot 행렬을 통해 달성되며, 동일한 위상 주기 \(\Delta N = 16.5\) ( \(L_0=96\) 및 \(\Seven\)에서 유도)를 가집니다.
		\item \textbf{정보 이론:} 일반화된 섀넌-야쿠셰프 법칙 \(W = K_{\mathrm{eff}} \cdot I_{\text{채널}} \cdot \eta\)은 \(E = mc^2\)의 직접적인 정보론적 유사체를 제공합니다.
		\item \textbf{기하학 및 우주론:} 회전 디스크에 대한 유효 계량은 \(d\ell^2 = dr^2 + r^2 K_{\mathrm{eff}}(r)^2 d\phi^2\)이며, 전체 우주 회전 주기 \(T_{\mathrm{rot}} \approx 2.3\times10^{14}\)년도 동일한 상수에서 따릅니다.
	\end{itemize}
	
	각 경우에 지배 방정식은 동일합니다:
	
	\begin{equation}
		\boxed{ \varepsilon = \kappa_c \, \alpha \, K_{\mathrm{eff}}^{-2/3} },
		\qquad \kappa_c = \frac{1}{3},\quad \beta = \frac{2}{3},
	\end{equation}
	
	여기서 \(\Keff\)는 분야에 따라 해석됩니다: 행동 엔트로피 대 키 엔트로피의 비율 (정보), 질량 비율 (입자 물리학), 대사 효율 (생물학), 응력 조정 (지구 물리학), 조정 깊이 (정수론) 또는 계량 계수 (기하학). 지수의 보편성은 이 모든 시스템이 YPSDC 프로토콜을 통해 작동하는 단일 조정 필드 \(\Psi_{MN}\)의 투영이기 때문에 발생합니다: 미리 분배된 사전이 짧은 키에 의해 활성화됩니다.
	
	따라서 YUCT는 자연 법칙의 명백한 다양성이 우리의 제한된 관점의 결과임을 드러냅니다. 기본 수준에서는 오직 하나의 법칙, 즉 조정의 법칙만이 존재합니다. 정보, 에너지, 기하학 및 순수 수학의 구조조차도 단일 조정 과정의 다른 측면입니다. 대수 순환 \(\Sodd + \Seven = 2\), \(\Sodd/\Seven = 3/2\), 및 \(\beta = 2/3\)는 이 통일된 사실의 견고한 구조를 제공하며, 50자리 이상의 규모에 걸친 실험적 확인은 이 통일성이 철학적 추측이 아니라 자연에서 측정 가능한 사실임을 보여줍니다.
	
	% ============================================================
	% 그림: 통일된 베타 (한국어 버전, 수정됨: gold 색상 정의됨)
	% ============================================================
	\begin{figure}[H]
		\centering
		\begin{tikzpicture}[
			node distance=1.2cm,
			dom/.style={rectangle, draw=blue!70, fill=blue!4, thick,
				minimum width=2.8cm, minimum height=0.7cm, rounded corners, align=center, font=\small},
			center/.style={rectangle, draw=gold!80!black, fill=gold!15, thick,
				minimum width=3.2cm, minimum height=1.2cm, rounded corners, align=center, font=\bfseries\large},
			arrow/.style={->, >=Stealth, thick, color=darkblue!70},
			label/.style={font=\scriptsize\itshape, align=center}
			]
			
			% 중앙 노드 – 보편적 법칙
			\node[center] (law) at (0,0) {\(\displaystyle \varepsilon = \kappa_c \alpha K_{\mathrm{eff}}^{-2/3}\)\\[2pt]
				\small\mdseries \(\beta = 2/3,\; \kappa_c = 1/3\)};
			
			% 분야 (원형 배치) — 한국어 텍스트
			\node[dom] (info) at (4.2, 2.8) {정보 이론\\부록 X};
			\node[dom] (particle) at (5.2, 0.0) {입자 물리학\\질량 사다리};
			\node[dom] (biology) at (4.2, -2.8) {생물학\\DNA, 대사};
			\node[dom] (geo) at (0.0, -4.5) {지구 물리학\\지진};
			\node[dom] (spiral) at (-4.2, -2.8) {나선형 구조\\14개 시스템};
			\node[dom] (number) at (-5.2, 0.0) {정수론\\소수};
			\node[dom] (quantum) at (-4.2, 2.8) {양자 역학\\Stern-Brocot};
			\node[dom] (cosmo) at (0.0, 4.5) {우주론\\회전, 인플레이션};
			
			% 법칙에서 분야로 화살표
			\draw[arrow] (law) -- (info);
			\draw[arrow] (law) -- (particle);
			\draw[arrow] (law) -- (biology);
			\draw[arrow] (law) -- (geo);
			\draw[arrow] (law) -- (spiral);
			\draw[arrow] (law) -- (number);
			\draw[arrow] (law) -- (quantum);
			\draw[arrow] (law) -- (cosmo);
			
			% 화살표 레이블 (한국어)
			\node[label, right] at (2.1, 1.6) {\(K_{\mathrm{eff}} = H(\mathcal A)/H(\mathcal K)\)};
			\node[label, right] at (3.0, 0.0) {\(K_{\mathrm{eff}} \propto m\)};
			\node[label, right] at (2.1, -1.6) {\(K_{\mathrm{eff}}\) = 복구 효율};
			\node[label, below] at (0.0, -2.5) {\(K_{\mathrm{eff}}\) = 응력 조정};
			\node[label, left] at (-2.1, -1.6) {\(K_{\mathrm{eff}}\) = 안정성 조건};
			\node[label, left] at (-3.0, 0.0) {\(K_{\mathrm{eff}} \propto p_n\)};
			\node[label, left] at (-2.1, 1.6) {\(K_{\mathrm{eff}}\) = 사전 크기};
			\node[label, above] at (0.0, 2.5) {\(K_{\mathrm{eff}}\) = 계량 계수};
			
			% 외부 프레임 및 제목
			\draw[gray!30, thick, rounded corners] (-6.9,-5.5) rectangle (6.9,5.8);
			\node[font=\small\bfseries, anchor=north west] at (-5.0,6.5) 
			{모든 분야에 걸친 보편적 조정 법칙};
			
		\end{tikzpicture}
		\caption{동일한 지수 \(\beta = 2/3\)가 정보 이론, 입자 물리학, 생물학, 지구 물리학, 유체 역학, 정수론, 양자 역학 및 우주론을 지배합니다. 각 분야에서 \(\Keff\)는 특정 해석을 취하지만, 함수 형태 \(\varepsilon = \kappa_c \alpha K_{\mathrm{eff}}^{-2/3}\)는 변하지 않습니다. 이 보편성은 단일 기본 조정 필드 \(\Psi_{MN}\)의 표시입니다.}
		\label{fig:unified-beta}
	\end{figure}
	
	\section{검증 가능한 예측}
	\label{sec:predictions}
	
	법칙 \(\eps = \kappac\,\alpha\,\Keff^{-2/3}\) with \(\kappac = 1/3\)에 기반하여, 우리는 다음과 같은 실험을 제안합니다:
	
	\begin{enumerate}
		\item \textbf{효소 촉매:} 효소 돌연변이 세트에 대해 \(\Keff\) (예: 컨포메이션 상태의 수로부터)와 촉매 효율 \(k_{\text{cat}}/K_M\)을 측정합니다. \(\log(k_{\text{cat}}/K_M)\) 대 \(\log \Keff\)의 기울기는 \(2/3\)여야 합니다. 더 자세한 생물학적 예측은 부록 D 및 E에 있습니다.
		\item \textbf{신경 코딩:} 신경 앙상블에서 자극-반응 상호 정보를 반응 엔트로피로 나눈 값으로 \(\Keff\)를 추정합니다. 판별 역치 (오류율)는 \(\Keff^{-2/3}\)로 스케일링되어야 합니다.
		\item \textbf{금융 시장:} 고주파 데이터를 사용하여 호가장 유동성과 매수-매도 스프레드로부터 \(\Keff\)를 계산합니다. 변동성 (수익률의 표준 편차)은 \(\sigma \propto \Keff^{-2/3}\)를 따를 것으로 예상됩니다.
		\item \textbf{우주 구조:} 유효 밀도 요동 \(\sigma_8\)은 초기 우주의 독립 영역 수에서 추정된 암흑 물질의 조정 효율과 상관되어야 합니다. 이는 CMB 관측과의 일관성 검사를 제공합니다 (본문 및 부록 G 참조).
		\item \textbf{초전도 회전 디스크 (부록 Ehrenfest):} 고온 초전도체로 만들어진 디스크의 경우, 붕괴되는 임계 각속도는 물질의 조정 효율 \(K_{\mathrm{eff,mat}}\)이 변하기 때문에 정상 상태와 초전도 상태 사이에서 달라져야 합니다. 이는 기하학의 물질 의존성에 대한 직접적인 실험실 테스트를 제공합니다.
	\end{enumerate}
	
	\section{한계 및 주의사항}
	\label{sec:limitations}
	
	보편적 스케일링 법칙을 적용할 때 주의가 필요합니다:
	\begin{itemize}
		\item \(\Keff\)의 정의는 사전과 행동 공간의 명확한 식별을 요구합니다; 많은 시스템에서 이는 사소하지 않습니다.
		\item 전인자 \(\kappac = 1/3\)는 삼원 구조에 대해 정밀합니다. YUCT에서 \(\alpha\)는 부록 Y, 섹션 Y.5.2에서 유도된 고정 상수입니다; 자유 매개변수가 아닙니다. 그러나 실제로 \(\Keff\)가 근사적으로만 추정될 때, 데이터에서 추론된 \(\alpha\)의 유효 값은 변할 수 있습니다. 이는 \(\alpha\)의 근본적인 변동이 아니라 \(\Keff\)의 더 정확한 추정의 필요성을 나타냅니다.
		\item 지수 \(2/3\)는 공간의 3차원성 (정리 2)과 연결 조건 (부록 Y)에서 엄밀히 유도되지만, \(\Keff\)에서 관측 가능한 양으로의 사영은 추가 가정을 필요로 할 수 있습니다.
		\item 표~\ref{tab:data}의 예는 오차가 직접적인 관심 대상이 아닌 시스템 (예: CMB)을 포함하므로, 정확한 사영이 필요합니다.
	\end{itemize}
	
	그럼에도 불구하고, 50자리 이상의 규모에 걸친 데이터의 수렴은 진정한 보편적 현상의 존재를 강력히 시사합니다.
	
	\section{결론}
	우리는 YUCT 내에서 보편적 오차 스케일링 법칙의 수정되고 수학적으로 일관된 공식을 제시했습니다. 지수 \(\beta = 2/3\)는 정밀하며 조정 공간의 3차원성 (본문 정리 2)과 계층 연결 조건 (부록 Y)에서 비롯되며, 전인자 \(\kappac = 1/3\)는 삼원 구조를 반영합니다. 실험값 \(0.67 \pm 0.03\)은 이 예측을 확인합니다. 이 법칙은 복잡계의 광범위한 멱법칙을 통일하고 검증 가능한 예측을 제공합니다. 미래 작업은 다양한 분야에서 \(\Keff\)의 정밀 측정에 초점을 맞추고 이론적 유도를 전인자 \(\alpha\)를 포함하도록 확장해야 합니다.
	
	\paragraph*{데이터 가용성}
	모든 데이터 수집 및 분석 스크립트는 \url{https://github.com/Alexey-Yakushev-YUCT/YPSDC}에서 제공됩니다.
	
	\paragraph*{버전 기록}
	\begin{itemize}
		\item v1.0: 초기 버전, 일관되지 않은 정의.
		\item v2.0: 정의 수정, 검증 예제 추가.
		\item v2.2: 원자 및 생물학적 규모로 확장.
		\item v3.0: 정리 2를 사용한 유도 수정, 실험값 \(0.67\)을 정밀 지수 \(2/3\)로 대체, 한계 명확화, 부록 P, Y, Ferra1.2에 대한 참조 추가.
		\item v3.1: 추가 명확화: 부록 Y에서 인자 2의 명시적 유도 추가, 오해의 소지가 있는 중성자 붕괴 예제 제거, 이징 모델 연관성을 독립적 확인으로 재구성, 작동화를 위해 부록 D, E, F에 대한 참조 추가, \(\alpha\)의 논의를 고정 상수로서의 지위를 반영하도록 수정.
		\item v3.2: 나선형 구조 및 소수를 포함한 분야 전반의 보편적 조정 섹션 추가; 부록 Ehrenfest 및 X와 연결; 관련 문헌 참조 추가.
		\item v3.3: 모든 분야를 명시적으로 통합하는 종합 하위 섹션 추가, 정보 이론, 입자 물리학, 생물학, 지구 물리학, 유체 역학, 정수론, 양자 역학 및 우주론에서 \(\beta = 2/3\)의 중심 역할을 보여주는 시각적 요약 그림 (그림~\ref{fig:unified-beta}) 추가.
	\end{itemize}
	
	\begin{thebibliography}{99}
		\bibitem{Yakushev2026}
		A. V. Yakushev, \emph{Yakushev Unified Coordination Theory (YUCT)}, Zenodo, \url{https://doi.org/10.5281/zenodo.18444599} (2026).
		\bibitem{AppendixC1}
		A. V. Yakushev, 부록 C1: 화학 결합 (HF–HI)에서 보편적 스케일링 법칙의 실험적 검증, \emph{YUCT} 내, Zenodo (2026).
		\bibitem{AppendixD}
		A. V. Yakushev, 부록 D: YUCT의 생물학적 예측 – 검증 가능한 가설, \emph{YUCT} 내, Zenodo (2026).
		\bibitem{AppendixE}
		A. V. Yakushev, 부록 E: 조정 유전학 – 분자 생물학의 YUCT 원리, \emph{YUCT} 내, Zenodo (2026).
		\bibitem{AppendixF}
		A. V. Yakushev, 부록 F: YUCT의 경제 응용, \emph{YUCT} 내, Zenodo (2026).
		\bibitem{AppendixG}
		A. V. Yakushev, 부록 G: Yakushev 통합 조정 이론 – 양자 중력 문제에 대한 근본적 해결, \emph{YUCT} 내, Zenodo (2026).
		\bibitem{AppendixP}
		A. V. Yakushev, 부록 P: 중력의 온도 의존성, \emph{YUCT} 내, Zenodo (2026).
		\bibitem{AppendixY}
		A. V. Yakushev, 부록 Y: YUCT의 수학적 기초 – 최초 원리로부터의 유도, \emph{YUCT} 내, Zenodo (2026).
		\bibitem{AppendixFerra}
		A. V. Yakushev, 부록 Ferra1.2: 질서 및 무질서 상에 대한 보편적 조정 상수, \emph{YUCT} 내, Zenodo (2026).
		\bibitem{AppendixPrimeN}
		A. V. Yakushev, 부록 PrimeN: YUCT 소수 조정 사다리, \emph{YUCT} 내, Zenodo (2026).
		\bibitem{AppendixX}
		A. V. Yakushev, 부록 X: 일반화된 섀넌 이론, 일정 오차 법칙 및 \(K_{\mathrm{eff}}\)-의존적 벨 부등식, \emph{YUCT} 내, Zenodo (2026).
		\bibitem{AppendixEhrenfest}
		A. V. Yakushev, 부록 Ehrenfest: 회전 디스크 역설의 조정 해석 및 비유클리드 기하학의 출현, \emph{YUCT} 내, Zenodo (2026).
		\bibitem{El-Showk2014}
		S. El-Showk et al., \emph{Phys. Rev. D} \textbf{90}, 105001 (2014).
		\bibitem{ammonites}
		D. K. Jacobs, Morphology and the evolution of ammonite shell coiling, \emph{Paleobiology} \textbf{16}, 283 (1990).
		\bibitem{watsoncrick}
		J. D. Watson and F. H. C. Crick, Molecular structure of nucleic acids, \emph{Nature} \textbf{171}, 737 (1953).
		\bibitem{pauling}
		L. Pauling and R. B. Corey, Configurations of polypeptide chains with favored orientations around single bonds, \emph{Proc. Natl. Acad. Sci. USA} \textbf{37}, 729 (1951).
		\bibitem{collagen}
		G. N. Ramachandran and G. Kartha, Structure of collagen, \emph{Nature} \textbf{176}, 593 (1955).
		\bibitem{phyllotaxis}
		P. H. Rivier and A. M. Turing, The chemical basis of morphogenesis, \emph{Phil. Trans. R. Soc. B} \textbf{237}, 37 (1952).
		\bibitem{cholesteric}
		P. G. de Gennes and J. Prost, \emph{The Physics of Liquid Crystals}, Oxford University Press, 2nd ed., 1993.
		\bibitem{helical-mof}
		B. Moulton and M. J. Zaworotko, From molecules to crystal engineering: supramolecular isomerism and polymorphism in coordination polymers, \emph{Chem. Rev.} \textbf{101}, 1629 (2001).
		\bibitem{ocean-eddies}
		A. R. Robinson (ed.), \emph{Eddies in Marine Science}, Springer, 1983.
		\bibitem{hurricanes}
		K. Emanuel, The physics of tropical cyclones, \emph{Annu. Rev. Fluid Mech.} \textbf{35}, 1 (2003).
		\bibitem{tornadoes}
		R. Davies-Jones, A review of supercell and tornado dynamics, \emph{Atmos. Res.} \textbf{158--159}, 274 (2015).
		\bibitem{nuclear-tracks}
		R. L. Fleischer, P. B. Price, and R. M. Walker, \emph{Nuclear Tracks in Solids: Principles and Applications}, University of California Press, 1975.
		\bibitem{binney-tremaine}
		J. Binney and S. Tremaine, \emph{Galactic Dynamics}, Princeton University Press, 2nd ed., 2008.
		\bibitem{ppdisk}
		S. M. Andrews and D. J. Wilner, The structure of protoplanetary disks, \emph{Annu. Rev. Astron. Astrophys.} \textbf{55}, 471 (2017).
		\bibitem{beck}
		R. Beck, Magnetic fields in galaxies, \emph{Space Sci. Rev.} \textbf{99}, 243 (2001).
		\bibitem{Tegmark2014}
		Max Tegmark, \emph{Our Mathematical Universe: My Quest for the Ultimate Nature of Reality}, Knopf, 2014.
	\end{thebibliography}
	
\end{document}